% Motor: XeLaTeX. Sesión 25 - material docente UTP.
% Estructura docente de la Sesión 24.
\documentclass[aspectratio=169,10pt,c]{beamer}
% Reutiliza sin cambios la identidad visual UTP validada en la Sesion 08.
% Reutiliza sin cambios la identidad visual UTP validada en la Sesion 05.
% La segunda ruta permite compartir los assets comunes sin duplicar binarios.
% !TEX root = slides.tex
% Preambulo comun: estilo UTP, colores, fuentes, plantillas de beamer y
% comandos de apoyo didactico. Se carga desde slides.tex tras \documentclass.

% La opción "c" de beamer reparte el espacio libre 1:1.5 (arriba:abajo), por lo
% que el contenido queda por encima del centro. Igualamos ambos rellenos.
\makeatletter
\define@key{beamerframe}{c}[true]{%
  \beamer@frametopskip=0pt plus 1fill\relax%
  \beamer@framebottomskip=0pt plus 1fill\relax%
  \beamer@frametopskipautobreak=0pt plus .5\paperheight\relax%
  \beamer@framebottomskipautobreak=0pt plus .5\paperheight\relax%
}
\makeatother

\usepackage{graphicx}
\usepackage{tikz}
\usepackage{xcolor}
\usepackage{iftex}

\ifPDFTeX
  \usepackage[utf8]{inputenc}
  \usepackage[T1]{fontenc}
  \usepackage{lmodern}          % tipografia escalable para el respaldo con pdflatex
  \renewcommand{\familydefault}{\sfdefault}
\else
  \usepackage{fontspec}
  \IfFontExistsTF{Arial}
    {\setsansfont{Arial}[BoldFont={Arial Bold}]}
    {\setsansfont{Helvetica}[BoldFont={Helvetica Bold}]}
\fi

\usepackage[spanish,es-nodecimaldot]{babel}
\usepackage{microtype}
\usepackage{booktabs}
\usepackage{array}
\usetikzlibrary{arrows.meta,positioning,fit,calc,shapes.geometric}

\graphicspath{{assets/}}

% ---------- Identidad visual UTP ----------
\definecolor{utpRojo}{HTML}{C42B35}
\definecolor{utpNegro}{HTML}{111111}
\definecolor{utpGris}{HTML}{5F5F5F}
\definecolor{utpGrisMedio}{HTML}{9A9A9A}
\definecolor{utpGrisClaro}{HTML}{F1F1F1}

\usetheme{default}
\setbeamercolor{normal text}{fg=utpNegro,bg=white}
\setbeamercolor{structure}{fg=utpRojo}
\setbeamercolor{frametitle}{fg=utpNegro,bg=white}
\setbeamercolor{title}{fg=utpNegro}
\setbeamercolor{subtitle}{fg=utpNegro}
\setbeamercolor{author}{fg=utpNegro}
\setbeamercolor{date}{fg=utpNegro}
\setbeamercolor{block title}{fg=white,bg=utpNegro}
\setbeamercolor{block body}{fg=utpNegro,bg=utpGrisClaro}
\setbeamercolor{alertblock title}{fg=white,bg=utpRojo}
\setbeamercolor{alertblock body}{fg=utpNegro,bg=utpRojo!7}
\setbeamercolor{exampleblock title}{fg=white,bg=utpGris}
\setbeamercolor{exampleblock body}{fg=utpNegro,bg=utpGris!8}
\setbeamercolor{terminal}{fg=white,bg=utpNegro}

\setbeamerfont{title}{series=\bfseries,size=\fontsize{19.5}{22.5}\selectfont}
\setbeamerfont{subtitle}{size=\fontsize{11.5}{13.5}\selectfont}
\setbeamerfont{date}{series=\bfseries,size=\fontsize{11.5}{13.5}\selectfont}
\setbeamerfont{author}{size=\fontsize{10.5}{12.5}\selectfont}
\setbeamerfont{frametitle}{series=\bfseries,size=\large}
\setbeamerfont{block title}{series=\bfseries,size=\small}
\setbeamerfont{block body}{size=\normalsize}

\setbeamertemplate{navigation symbols}{}
\setbeamertemplate{headline}{}
\setbeamertemplate{blocks}[default]
\setbeamertemplate{itemize item}{\color{utpRojo}\small\raise.15ex\hbox{$\blacktriangleright$}}
\setbeamertemplate{itemize subitem}{\color{utpGris}\scriptsize$\bullet$}
\setbeamertemplate{enumerate item}{\color{utpRojo}\insertenumlabel.}
\setbeamertemplate{background canvas}{\color{white}\rule{\paperwidth}{\paperheight}}

% Título de diapositiva: negrita negra a la izquierda, logo UTP a la derecha
\setbeamertemplate{frametitle}{%
  \begin{tikzpicture}[remember picture,overlay]
    \node[anchor=north east,inner sep=0pt]
      at ([xshift=-7mm,yshift=-5.5mm]current page.north east)
      {\includegraphics[width=17mm]{utp-logo.png}};
  \end{tikzpicture}%
  \vspace{1.2mm}%
  \begin{minipage}[t]{.80\textwidth}\raggedright\insertframetitle\end{minipage}\par
  \vspace{1.1mm}{\color{utpRojo}\rule{\textwidth}{1.1pt}}
}

% Solo numeración discreta; la plantilla UTP no lleva pie de página
\setbeamertemplate{footline}{%
  \hfill{\color{utpGrisMedio}\scriptsize\insertframenumber}\hspace{6mm}\vspace{2.5mm}%
}

% Portada UTP
\setbeamertemplate{title page}{%
  \begin{tikzpicture}[remember picture,overlay]
    \node[anchor=center,inner sep=0pt] at (current page.center) {%
      \includegraphics[width=\paperwidth,height=\paperheight]{utp-cover-bg.jpeg}%
    };
    \node[anchor=north,align=center,text=utpNegro,inner sep=0pt,
      text width=145mm,font=\usebeamerfont{title}]
      at ([yshift=-25.2mm]current page.north) {\inserttitle};
    \node[anchor=north,align=center,text=utpNegro,inner sep=0pt,
      text width=142mm,font=\usebeamerfont{subtitle}]
      at ([yshift=-38.2mm]current page.north) {\insertsubtitle};
    \node[anchor=north,align=center,text=utpNegro,inner sep=0pt,
      font=\usebeamerfont{date}]
      at ([yshift=-50.4mm]current page.north) {\insertdate};
    \node[anchor=north,align=center,text=utpNegro,inner sep=0pt,
      font=\usebeamerfont{author}]
      at ([yshift=-58.4mm]current page.north) {\insertauthor};
  \end{tikzpicture}%
}

% ---------- Comandos de apoyo didáctico ----------
\newcommand{\concepto}[2]{%
  \begin{block}{#1}\large #2\end{block}}
\newcommand{\pregunta}[1]{%
  \begin{alertblock}{Decide y explica}#1\end{alertblock}}
\newcommand{\analogia}[1]{%
  \begin{exampleblock}{En palabras simples}\small #1\end{exampleblock}}
\newcommand{\enlarealidad}[1]{%
  \begin{block}{Cómo se ve en la realidad}\small #1\end{block}}
\newcommand{\ojo}[1]{%
  \begin{alertblock}{Ojo con esto}\small #1\end{alertblock}}
\newcommand{\apoyo}{{\normalfont\footnotesize\color{utpGris}~· material de apoyo}}
\newcommand{\marca}[1]{\tikz[baseline=(m.base)]{%
  \node[rounded corners=1.5pt,fill=utpGrisClaro,draw=utpGrisMedio,inner xsep=3pt,inner ysep=1.5pt,
        font=\footnotesize] (m) {#1};}}
\newenvironment{terminal}{%
  \begin{beamercolorbox}[wd=\linewidth,sep=5pt,rounded=true]{terminal}%
  \ttfamily\small\ignorespaces}{%
  \end{beamercolorbox}}

\tikzset{
  caja/.style={draw=utpNegro, very thick, rounded corners=1.5pt, fill=white,
    align=center, minimum height=0.72cm, inner sep=5pt, font=\small},
  cajaazul/.style={caja, draw=utpGris, fill=white},
  cajaverde/.style={caja, draw=utpNegro, fill=utpGrisClaro},
  cajanaranja/.style={caja, draw=utpRojo!75, fill=utpRojo!7},
  flecha/.style={-{Stealth[length=2.3mm]}, very thick, draw=utpNegro},
  flechavuelta/.style={{Stealth[length=2.3mm]}-{Stealth[length=2.3mm]}, very thick, draw=utpNegro}
}

\graphicspath{{../../Sesion05/slides/assets/}}



\graphicspath{{assets/}{../../Sesion05/slides/assets/}}
\usepackage{colortbl}
\usepackage{listings}
\lstset{basicstyle=\ttfamily\footnotesize,backgroundcolor=\color{utpGrisClaro},
  frame=none,aboveskip=4pt,belowskip=4pt,showstringspaces=false,
  columns=fullflexible,keepspaces=true,breaklines=true,
  keywordstyle=\color{utpRojo}\bfseries,
  morekeywords={SELECT,FROM,GROUP,BY,ORDER,DESC,AS,SUM,ROUND,import},
  xleftmargin=5pt,xrightmargin=5pt,framesep=5pt}

\title{Integración de Servicios Cloud GCP}
\subtitle{Supervisión, alta disponibilidad y balanceo de carga}
\date{Semana 13 - Sesión 25}
\author{Rel Guzman}

\newcommand{\separador}[2]{%
\begin{frame}[plain]
\vfill
\begin{center}
{\fontsize{29}{34}\selectfont\bfseries #1}\par
\vspace{2mm}
{\small #2}\par
\vspace{4mm}
{\color{utpRojo}\rule{8.5cm}{1.2pt}}
\end{center}
\vfill
\end{frame}}

% Íconos simples para los diagramas.
\tikzset{
 caja/.style={draw=utpRojo,thick,rounded corners=4pt,align=center},
 vm/.style={draw=utpRojo,thick,rounded corners=3pt,fill=white,minimum width=1.6cm,minimum height=0.75cm,align=center,font=\scriptsize},
 vmmal/.style={vm,draw=utpGrisMedio,fill=utpGrisClaro,text=utpGris},
 zona/.style={draw=utpGrisMedio,dashed,rounded corners=4pt},
 flecha/.style={-latex,thick},
}
\newcommand{\cruz}[1]{\draw[utpRojo,line width=2pt] ($(#1)+(-0.25,-0.25)$) -- ($(#1)+(0.25,0.25)$) ($(#1)+(-0.25,0.25)$) -- ($(#1)+(0.25,-0.25)$);}

\begin{document}

{\setbeamertemplate{footline}{}\maketitle}

% 2
\begin{frame}[fragile]{Revisión: BigQuery sin BigQuery}
\begin{columns}[T,onlytextwidth]
\column{.44\textwidth}
\begin{block}{Qué harás}
\small Las consultas del lab de la sesión anterior, ahora con \textbf{Python} y \textbf{DuckDB}: un motor SQL que corre dentro de Python y, como BigQuery, lee por columnas. No necesitas cuenta de Google Cloud.
\end{block}
\begin{block}{De dónde salen los datos}
\small No se descargan: la primera celda del notebook ejecuta \texttt{generar\_ventas.py}, que crea \texttt{ventas.csv} con \textbf{1.2 millones} de pedidos de Pedidos Express (2023--2025). Todos obtienen los mismos datos.
\end{block}
\column{.53\textwidth}
\begin{lstlisting}
import duckdb

duckdb.sql("""
  SELECT categoria,
    ROUND(SUM(monto), 2) AS total_soles
  FROM ventas
  GROUP BY categoria
  ORDER BY total_soles DESC
""")
\end{lstlisting}
\par\vspace{2mm}
\begin{exampleblock}{Cómo empezar}
\scriptsize 1. En Colab: \textbf{Archivo} $\rightarrow$ \textbf{Abrir cuaderno} $\rightarrow$ \textbf{GitHub}.\par
2. Pega el repositorio del curso:\par
\hspace*{3mm}\texttt{github.com/rgap/}\par
\hspace*{3mm}\texttt{rgap-Integracion-de-Servicios-Cloud-Gcp-Agosto2026}\par
3. Elige \texttt{Sesion25/practica/practica\_bigquery\_python.ipynb}, ejecuta las celdas en orden y resuelve los 5 retos del final.
\end{exampleblock}
\end{columns}
\end{frame}

% 3
\begin{frame}{Revisión: compara tus resultados}
\begin{columns}[T,onlytextwidth]
\column{.52\textwidth}
\renewcommand{\arraystretch}{1.35}
\small
\begin{tabular}{@{}>{\raggedright\arraybackslash}p{4.4cm}p{2.3cm}@{}}
\toprule
\textbf{Pregunta del notebook} & \textbf{Tu resultado}\\
\midrule
¿Cuántas filas tiene la tabla? & \\
4.3 ¿Qué categoría tiene más pedidos? & \\
4.5 ¿Qué categoría vende más dinero? & \\
Paso 5: ¿cuánto ocupa el CSV? ¿Y el Parquet? & \\
\bottomrule
\end{tabular}
\column{.44\textwidth}
\begin{alertblock}{Responde}
\small
1. ¿La categoría con más pedidos es la que más dinero vende? ¿Por qué?\par\medskip
2. ¿Por qué la consulta sobre Parquet termina antes que la consulta sobre CSV?\par\medskip
3. ¿Qué tienen en común DuckDB y BigQuery?
\end{alertblock}
\end{columns}
\end{frame}

% 4
\begin{frame}{Revisión: la pregunta de hoy}
\begin{center}
\begin{tikzpicture}[font=\small]
  \foreach \i in {0,1,2} \node[caja,fill=utpRojo!12,minimum width=1.6cm] (c\i) at (0,1.1-\i*1.1) {cliente};
  \node[vm,minimum width=3.2cm,minimum height=1.2cm,font=\small] (w) at (5.2,0) {\textbf{VM} \texttt{web-1}\\\footnotesize zona \texttt{us-central1-a}};
  \foreach \i in {0,1,2} \draw[flecha] (c\i.east) -- (w.west);
  \node[caja,fill=utpRojo!10,text width=3.9cm,font=\footnotesize] at (9.6,0) {Viernes, 20:00. Hora de la cena.\\Llegan 10 veces más pedidos.\\\textbf{La página no carga.}};
\end{tikzpicture}
\end{center}
\vspace{2mm}
\begin{columns}[T,onlytextwidth]
\column{.32\textwidth}
\begin{block}{1. Supervisar}
\small ¿Cómo te enteras \textbf{antes} de que llamen los clientes?
\end{block}
\column{.32\textwidth}
\begin{block}{2. Alta disponibilidad}
\small Si \texttt{web-1} se apaga, ¿qué sigue respondiendo?
\end{block}
\column{.32\textwidth}
\begin{block}{3. Balanceo de carga}
\small Si pones 3 VMs, ¿cómo se reparten los clientes entre ellas?
\end{block}
\end{columns}
\end{frame}

% 5
\begin{frame}[c]
\begin{block}{Logro de la sesión}
Al finalizar, podrás explicar qué es \textbf{supervisar recursos} y qué ofrece \textbf{Google Cloud Observability}, crear en \textbf{Cloud Monitoring} un panel, una alerta y una verificación de tiempo de actividad, y explicar cómo la \textbf{alta disponibilidad}, el \textbf{balanceador de carga} y \textbf{Cloud CDN} mantienen una aplicación en línea.
\end{block}
\end{frame}

% 6
\begin{frame}{Saberes previos}
\begin{itemize}\setlength\itemsep{11pt}
  \item Si una página web tarda 10 segundos en cargar, ¿qué podría estar pasando en el servidor?
  \item Creaste una VM en la zona \texttt{us-central1-a}. ¿Qué le pasa a tu página si esa zona tiene una falla?
  \item Tienes tres VMs con la misma página y cada una tiene su propia IP externa. ¿Qué dirección les das a los clientes?
\end{itemize}
\end{frame}

% 7
\begin{frame}[c]{Utilidad de esta sesión}
\begin{center}
\begin{minipage}{.84\textwidth}
{\large Toda aplicación en producción falla alguna vez: se llena un disco, la CPU llega al 100\,\%, una VM se reinicia. Si sabes \textbf{supervisar}, te enteras por una alerta y no por un cliente. Si sabes repartir la aplicación entre varias VMs detrás de un \textbf{balanceador}, una falla no deja a nadie sin servicio.}
\end{minipage}
\end{center}
\end{frame}

% 8
\separador{Parte 1: supervisar recursos}{Saber qué pasa en tus servidores antes de que te lo digan los clientes}

% 9
\begin{frame}{¿Qué es supervisar un recurso?}
\begin{block}{Definición}
\small Medir de forma continua el estado de los recursos (VMs, bases de datos, aplicaciones), guardar esas medidas y \textbf{avisar} cuando algo sale de lo normal.
\end{block}
\vspace{1mm}
\begin{center}
\begin{tikzpicture}[font=\small,
 p/.style={draw=utpRojo,thick,rounded corners=5pt,fill=white,text width=3.9cm,minimum height=3.3cm,align=center,anchor=north}]
  \node[p,fill=utpRojo!10] at (0,0) {{\Large\bfseries Métricas}\\[2pt]\footnotesize Números que cambian en el tiempo\\[8pt]\scriptsize\ttfamily web-1 · uso de CPU\\20:00 \quad 35 \%\\20:01 \quad 62 \%\\20:02 \quad 95 \%};
  \node[p] at (4.6,0) {{\Large\bfseries Registros}\\[2pt]\footnotesize Líneas de texto con lo que pasó (\textit{logs})\\[8pt]\scriptsize\ttfamily\raggedright 20:02:14 GET /menu 200\\20:02:15 GET /pagar 500\\20:02:15 ERROR: no se pudo conectar a la base};
  \node[p,fill=utpRojo!10] at (9.2,0) {{\Large\bfseries Alertas}\\[2pt]\footnotesize Reglas que avisan a una persona\\[8pt]\small Si el uso de CPU pasa de \textbf{80\,\%} durante \textbf{5 min}, envía un correo al equipo};
\end{tikzpicture}
\end{center}
\begin{alertblock}{Cada pieza responde una pregunta}
\small Métricas: \textbf{¿cuánto?} \quad Registros: \textbf{¿qué pasó exactamente?} \quad Alertas: \textbf{¿a quién aviso?}
\end{alertblock}
\end{frame}

% 10
\begin{frame}{Las métricas que más se revisan}
\renewcommand{\arraystretch}{1.05}
\centering\small
\begin{tabular}{@{}p{2.6cm}>{\raggedright\arraybackslash}p{4.6cm}>{\raggedright\arraybackslash}p{5.4cm}@{}}
\toprule
\textbf{Métrica} & \textbf{Qué mide} & \textbf{Señal de problema}\\
\midrule
\textbf{Uso de CPU} & Qué porcentaje del procesador se usa & Se queda sobre el 80\,\% por varios minutos\\
\textbf{Memoria} & Cuánta RAM ocupa la aplicación & Llega cerca del 100\,\%: la aplicación se cierra\\
\textbf{Disco} & Espacio usado del disco & Pasa del 90\,\%: pronto no se podrá guardar nada\\
\textbf{Red} & Bytes que entran y salen de la VM & Cae a 0 en horario de atención\\
\textbf{Latencia} & Cuánto tarda en llegar una respuesta & Sube de 200 ms a 3 s\\
\textbf{Errores} & Respuestas con código 5xx (falla del servidor) & Aparecen muchas de golpe\\
\bottomrule
\end{tabular}
\par\vspace{2mm}
\begin{alertblock}{Dato útil}
\small Una VM envía sola su uso de CPU, red y lecturas de disco. Para ver la \textbf{memoria} y el \textbf{espacio de disco} se instala en la VM el \textbf{Agente de operaciones} (\textit{Ops Agent}).
\end{alertblock}
\end{frame}

% 11
\begin{frame}{Leer una gráfica de uso de CPU}
\begin{center}
\begin{tikzpicture}[font=\scriptsize,x=0.46cm,y=0.032cm]
  % Ejes: horas del día (0 a 24) y porcentaje (0 a 100).
  \foreach \y in {0,20,40,60,80,100} {
    \draw[utpGrisClaro] (0,\y) -- (24,\y);
    \node[anchor=east,color=utpGris] at (0,\y) {\y\,\%};
  }
  \foreach \x in {0,4,8,12,16,20,24} \node[anchor=north,color=utpGris] at (\x,0) {\x:00};
  \draw[utpGrisMedio] (0,0) -- (24,0);
  \draw[utpNegro,thick,dashed] (0,80) -- (24,80);
  \node[anchor=south west,font=\scriptsize\bfseries,fill=white,inner sep=1pt] at (0.2,81) {umbral de la alerta: 80\,\%};
  \draw[utpRojo,very thick,line join=round]
    (0,8) -- (2,5) -- (4,5) -- (6,8) -- (8,20) -- (10,30) -- (12,62) -- (13,70) -- (14,48)
    -- (16,34) -- (18,52) -- (19,82) -- (20,95) -- (21,90) -- (22,50) -- (23,24) -- (24,12);
  \node[caja,fill=white,font=\scriptsize,anchor=east] (pc) at (17.6,104) {cena: 95\,\%, sobre el umbral};
  \draw[-latex] (pc.east) -- (19.8,96);
  \node[caja,fill=white,font=\scriptsize] (pa) at (9,60) {almuerzo: 70\,\%};
  \draw[-latex] (pa.east) -- (12.8,71);
  \node[caja,fill=white,font=\scriptsize] (pm) at (4,30) {madrugada: 5\,\%};
  \draw[-latex] (pm.south) -- (3,6);
\end{tikzpicture}
\end{center}
\vspace{1mm}
\begin{columns}[T,onlytextwidth]
\column{.48\textwidth}
\begin{block}{Lo normal}
\small La CPU sube en hora punta y baja de madrugada. Un pico corto no es una falla.
\end{block}
\column{.48\textwidth}
\begin{block}{Lo raro}
\small Que se quede alta mucho tiempo, que esté alta sin clientes o que la línea desaparezca (la VM no envía datos).
\end{block}
\end{columns}
\end{frame}

% 12
\begin{frame}{Problemas que ocurren seguido}
\renewcommand{\arraystretch}{1.3}
\centering\small
\begin{tabular}{@{}>{\raggedright\arraybackslash}p{3.7cm}>{\raggedright\arraybackslash}p{4.4cm}>{\raggedright\arraybackslash}p{4.6cm}@{}}
\toprule
\textbf{Qué pasa} & \textbf{Qué ves} & \textbf{Qué se hace}\\
\midrule
Hora punta: muchos más pedidos & CPU al 95\,\% y latencia alta & Agregar VMs (Parte 3)\\
Los registros llenan el disco & Disco al 100\,\% y errores al guardar & Borrar registros viejos o agrandar el disco\\
La VM se detuvo o se reinició & La verificación de tiempo de actividad falla; la gráfica se corta & Revisar los registros y tener otra VM (Parte 2)\\
Se subió una versión con un error & Errores 5xx justo después de la hora del cambio & Volver a la versión anterior\\
Un proceso quedó trabado & CPU al 100\,\% a las 4 a.\,m., sin clientes & Revisar el proceso y reiniciarlo\\
\bottomrule
\end{tabular}
\par\vspace{3mm}
\begin{alertblock}{El orden de trabajo}
\small La \textbf{alerta} avisa $\rightarrow$ la \textbf{métrica} dice desde cuándo y cuánto $\rightarrow$ los \textbf{registros} dicen por qué.
\end{alertblock}
\end{frame}

% 13
\begin{frame}{Google Cloud Observability}
\begin{block}{Definición}
\small Conjunto de servicios de Google Cloud para supervisar tus recursos. Ya viene activado en cada proyecto: las VMs, Cloud SQL y los buckets envían métricas y registros sin que instales nada.
\end{block}
\begin{center}
\begin{tikzpicture}[font=\small]
  \node[caja,fill=white,minimum width=2.6cm] (r1) at (0,1.2) {VMs};
  \node[caja,fill=white,minimum width=2.6cm] (r2) at (0,0.35) {Cloud SQL};
  \node[caja,fill=white,minimum width=2.6cm] (r3) at (0,-0.5) {Cloud Storage};
  \node[caja,fill=utpRojo!25,text width=5.4cm] (m) at (5.6,1.05) {\textbf{Cloud Monitoring}\\\footnotesize métricas, paneles, alertas y verificaciones\\de tiempo de actividad};
  \node[caja,fill=utpRojo!12,text width=5.4cm] (l) at (5.6,-0.45) {\textbf{Cloud Logging}\\\footnotesize registros; se buscan en el \textbf{Explorador de registros}};
  \foreach \r in {r1,r2,r3} { \draw[flecha] (\r.east) -- (m.west); \draw[flecha] (\r.east) -- (l.west); }
  \node[caja,fill=utpRojo!10,minimum width=2.0cm] (c) at (10.9,1.05) {correo\\\footnotesize al equipo};
  \draw[flecha] (m.east) -- node[above,font=\scriptsize]{alerta} (c.west);
\end{tikzpicture}
\end{center}
\renewcommand{\arraystretch}{1.15}
\centering\footnotesize
\begin{tabular}{@{}>{\raggedright\arraybackslash}p{6.2cm}>{\raggedright\arraybackslash}p{6.2cm}@{}}
\toprule
\textbf{Cloud Monitoring} {\color{utpGris}(lo usas en el lab)} & \textbf{Cloud Logging}\\
\midrule
¿Cuánta CPU usó \texttt{web-1} entre 19:00 y 21:00? & ¿Qué error dio la página \texttt{/pagar} a las 20:02?\\
\bottomrule
\end{tabular}
\end{frame}

% 14
\begin{frame}{Cloud Monitoring: lo que harás en el lab}
\begin{center}
\begin{tikzpicture}[font=\small,
 h/.style={draw=utpRojo,thick,rounded corners=5pt,fill=white,text width=5.6cm,minimum height=1.55cm,align=left,anchor=north west}]
  \node[h,fill=utpRojo!10] at (0,0) {\textbf{Panel} (\textit{dashboard})\\\footnotesize Una página con gráficos que tú eliges. En el lab: \texttt{My dashboard} con un gráfico de líneas del uso de CPU.};
  \node[h] at (6.3,0) {\textbf{Explorador de métricas}\\\footnotesize Ver una métrica al momento, sin guardarla en un panel.};
  \node[h] at (0,-1.85) {\textbf{Política de alertas}\\\footnotesize Una regla: métrica + umbral + a quién avisar. En el lab: \texttt{My Alert Policy}.};
  \node[h,fill=utpRojo!10] at (6.3,-1.85) {\textbf{Grupo}\\\footnotesize Junta recursos por su nombre para verlos juntos. En el lab: las VMs cuyo nombre contiene \texttt{nginx}.};
  \node[h,fill=utpRojo!10,text width=11.9cm] at (0,-3.7) {\textbf{Verificación de tiempo de actividad} (\textit{uptime check})\\\footnotesize Google pide tu página por HTTP cada minuto desde varias partes del mundo. Si no responde, te avisa. En el lab: \texttt{My Uptime check} sobre el grupo.};
\end{tikzpicture}
\end{center}
\vspace{-1mm}
{\footnotesize\color{utpGris} El lab ya trae tres VMs con un servidor web: \texttt{nginxstack-1}, \texttt{nginxstack-2} y \texttt{nginxstack-3}.}
\end{frame}

% 15
\begin{frame}{Cómo funciona una política de alertas}
\begin{center}
\begin{tikzpicture}[font=\small,
 paso/.style={caja,fill=white,text width=2.55cm,minimum height=1.9cm}]
  \node[paso,fill=utpRojo!10] (a) at (0,0) {\textbf{Métrica}\\[3pt]\footnotesize Instancia de VM\\$\rightarrow$ Uso de CPU};
  \node[paso] (b) at (3.4,0) {\textbf{Condición}\\[3pt]\footnotesize Por encima de \textbf{20\,\%}\\ventana de 1 min};
  \node[paso,fill=utpRojo!10] (c) at (6.8,0) {\textbf{Incidente}\\[3pt]\footnotesize se abre cuando la condición se cumple};
  \node[paso] (d) at (10.2,0) {\textbf{Canal de notificación}\\[3pt]\footnotesize correo electrónico};
  \draw[flecha] (a) -- (b); \draw[flecha] (b) -- (c); \draw[flecha] (c) -- (d);
\end{tikzpicture}
\end{center}
\vspace{2mm}
\begin{columns}[T,onlytextwidth]
\column{.48\textwidth}
\begin{block}{En el lab}
\small La política tiene \textbf{dos condiciones} y se dispara solo si se cumplen todas (\textit{All conditions are met}). El umbral es 20\,\% para que la veas dispararse; en producción se usa algo como 80\,\%.
\end{block}
\column{.48\textwidth}
\begin{block}{Elegir el umbral}
\small Muy bajo: llegan correos todo el día y el equipo deja de leerlos. Muy alto: avisa cuando los clientes ya se fueron.
\end{block}
\end{columns}
\par\vspace{2mm}
\begin{alertblock}{Al terminar el lab}
\small Desactiva la política (tarea 6). Las alertas siguen activas un tiempo después de borrar el proyecto.
\end{alertblock}
\end{frame}

% 16
\separador{Parte 2: alta disponibilidad}{Que la aplicación siga en línea aunque falle una pieza}

% 17
\begin{frame}{¿Qué es la disponibilidad?}
\begin{block}{Disponibilidad}
\small El porcentaje del tiempo en que el servicio responde. Un año tiene $365 \times 24 = 8\,760$ horas.
\end{block}
\renewcommand{\arraystretch}{1.3}
\centering\small
\begin{tabular}{@{}lrr>{\raggedright\arraybackslash}p{5.2cm}@{}}
\toprule
\textbf{Disponibilidad} & \textbf{Caída al año} & \textbf{Caída al mes} & \textbf{Ejemplo}\\
\midrule
99\,\% & 87.6 horas & 7.3 horas & Una VM sola que a veces se reinicia\\
99.9\,\% & 8.8 horas & 44 minutos & Varias VMs en una zona\\
99.99\,\% & 53 minutos & 4.4 minutos & Varias VMs en varias zonas\\
\bottomrule
\end{tabular}
\par\vspace{2mm}
{\scriptsize\color{utpGris} Cálculo: 1\,\% de 8\,760 horas = 87.6 horas. Los ejemplos son orientativos.}
\par\vspace{2mm}
\begin{alertblock}{Alta disponibilidad}
\small Diseñar el sistema para que la falla de \textbf{una pieza} (una VM, un disco, una zona) no apague el servicio.
\end{alertblock}
\end{frame}

% 18
\begin{frame}{Punto único de falla}
\begin{columns}[T,onlytextwidth]
\column{.42\textwidth}
\centering{\color{utpRojo}\bfseries Una VM en una zona}\par\vspace{2mm}
\begin{tikzpicture}[font=\small]
  \draw[zona] (-1.5,-0.8) rectangle (1.5,0.8);
  \node[font=\scriptsize,color=utpGris,anchor=north west] at (-1.5,0.8) {us-central1-a};
  \node[vmmal] (w) at (0,-0.1) {\texttt{web-1}};
  \cruz{w}
  \node[caja,fill=utpRojo!10,text width=3.6cm,font=\footnotesize] at (0,-1.8) {Falla \texttt{web-1} o la zona:\\\textbf{0 servidores, página caída}};
\end{tikzpicture}
\column{.56\textwidth}
\centering{\color{utpRojo}\bfseries Tres VMs en tres zonas}\par\vspace{2mm}
\begin{tikzpicture}[font=\small]
  \foreach \z/\n/\x in {a/1/0,b/2/2.3,c/3/4.6} {
    \draw[zona] (\x-1.05,-0.8) rectangle (\x+1.05,0.8);
    \node[font=\scriptsize,color=utpGris,anchor=north west] at (\x-1.05,0.8) {us-central1-\z};
  }
  \node[vmmal] (w1) at (0,-0.1) {\texttt{web-1}};
  \cruz{w1}
  \node[vm] at (2.3,-0.1) {\texttt{web-2}};
  \node[vm] at (4.6,-0.1) {\texttt{web-3}};
  \node[caja,fill=utpRojo!10,text width=5.4cm,font=\footnotesize] at (2.3,-1.8) {Falla la zona \texttt{a}:\\\textbf{quedan 2 servidores, la página sigue}};
\end{tikzpicture}
\end{columns}
\par\vspace{4mm}
\begin{block}{Punto único de falla}
\small Una pieza que, si falla, apaga todo el sistema. Para quitarlo: \textbf{duplicar} la pieza y poner las copias en \textbf{zonas distintas} de la región.
\end{block}
\end{frame}

% 19
\begin{frame}{Alta disponibilidad en Google Cloud}
\renewcommand{\arraystretch}{1.35}
\centering\small
\begin{tabular}{@{}p{2.6cm}>{\raggedright\arraybackslash}p{3.8cm}>{\raggedright\arraybackslash}p{6.2cm}@{}}
\toprule
\textbf{Pieza} & \textbf{Sin alta disponibilidad} & \textbf{Con alta disponibilidad}\\
\midrule
\textbf{Servidores web} & Una VM en una zona & \textbf{Grupo de instancias administrado}: varias VMs iguales, creadas desde una plantilla, en varias zonas. Si una VM se detiene, el grupo la vuelve a crear.\\
\textbf{Base de datos} & Cloud SQL en una zona & \textbf{Cloud SQL con alta disponibilidad}: una copia en espera en otra zona. Si la principal falla, la copia toma su lugar sola.\\
\textbf{Archivos} & Un disco en una VM & \textbf{Cloud Storage}: guarda varias copias de cada archivo sin que hagas nada.\\
\bottomrule
\end{tabular}
\par\vspace{3mm}
\begin{columns}[T,onlytextwidth]
\column{.48\textwidth}
\begin{block}{La regla}
\small Más de una copia, en zonas distintas, y algo que detecte la falla y cambie solo.
\end{block}
\column{.48\textwidth}
\begin{alertblock}{El costo}
\small Tres VMs cuestan tres veces una. Se elige según cuánto le cuesta al negocio una hora caída.
\end{alertblock}
\end{columns}
\end{frame}

% 20
\separador{Parte 3: balanceo de carga y CDN}{Una sola dirección para muchos servidores}

% 21
\begin{frame}{El balanceador de carga}
\begin{center}
\begin{tikzpicture}[font=\small]
  \foreach \i in {0,...,3} \node[caja,fill=utpRojo!12,minimum width=1.5cm,font=\footnotesize] (c\i) at (0,1.65-\i*1.1) {cliente};
  \node[caja,fill=utpRojo!30,text width=3.2cm,minimum height=1.4cm] (lb) at (4.3,0) {\textbf{Balanceador}\\\footnotesize una sola IP pública\\\texttt{34.120.10.5}};
  \foreach \i in {0,...,3} \draw[flecha] (c\i.east) -- (lb.west);
  \foreach \n/\y in {1/1.2,2/0,3/-1.2} {
    \node[vm] (w\n) at (9.0,\y) {\texttt{web-\n}};
    \draw[flecha] (lb.east) -- node[pos=0.68,above,sloped,font=\scriptsize]{100 por seg.} (w\n.west);
  }
  \node[font=\scriptsize,color=utpGris,align=center] at (4.3,-1.35) {300 solicitudes\\por segundo};
\end{tikzpicture}
\end{center}
\vspace{-2mm}
\begin{block}{Balanceador de carga}
\small Servicio que recibe todas las solicitudes en una sola dirección IP y las \textbf{reparte} entre varios servidores. Los clientes no saben cuántas VMs hay detrás.
\end{block}
\begin{columns}[T,onlytextwidth]
\column{.48\textwidth}
\begin{exampleblock}{Sin balanceador}
\small Tres IPs distintas. Si publicas la de \texttt{web-1}, las otras dos no reciben a nadie.
\end{exampleblock}
\column{.48\textwidth}
\begin{exampleblock}{Con balanceador}
\small Una IP. Cada VM atiende una parte; si agregas una cuarta, el balanceador la usa.
\end{exampleblock}
\end{columns}
\end{frame}

% 22
\begin{frame}{Verificación de estado: no enviar tráfico a una VM caída}
\begin{columns}[c,onlytextwidth]
\column{.44\textwidth}
\centering
\begin{tikzpicture}[font=\small]
  \node[caja,fill=utpRojo!30,minimum width=2.1cm,minimum height=1.1cm] (lb) at (0,0) {\textbf{Balanceador}};
  \node[vm] (w1) at (4.0,1.2) {\texttt{web-1}};
  \node[vmmal] (w2) at (4.0,0) {\texttt{web-2}};
  \node[vm] (w3) at (4.0,-1.2) {\texttt{web-3}};
  \draw[flecha] (lb.east) -- node[pos=0.62,above,sloped,font=\scriptsize]{150 por seg.} (w1.west);
  \draw[thick,dashed,utpGrisMedio] (lb.east) -- node[pos=0.55,above,font=\scriptsize,text=utpGris]{0} (w2.west);
  \draw[flecha] (lb.east) -- node[pos=0.62,below,sloped,font=\scriptsize]{150 por seg.} (w3.west);
  \cruz{w2}
\end{tikzpicture}
\column{.54\textwidth}
\renewcommand{\arraystretch}{1.2}
\scriptsize
\begin{tabular}{@{}lll@{}}
\toprule
\textbf{Hora} & \textbf{Pregunta a \texttt{web-2}} & \textbf{Estado}\\
\midrule
20:00:00 & \texttt{GET /} $\rightarrow$ 200 OK & sana\\
20:00:05 & \texttt{GET /} $\rightarrow$ sin respuesta & sana (1 falla)\\
\rowcolor{utpRojo!20} 20:00:10 & \texttt{GET /} $\rightarrow$ sin respuesta & \textbf{no sana}\\
\bottomrule
\end{tabular}
\par\vspace{1mm}
{\scriptsize\color{utpGris} Valores por defecto: pregunta cada 5 s; 2 fallas seguidas la marcan como no sana.}
\end{columns}
\par\vspace{3mm}
\begin{block}{Verificación de estado (\textit{health check})}
\small El balanceador pregunta a cada VM cada pocos segundos si responde. Si no responde, deja de enviarle clientes; cuando vuelve a responder, la usa otra vez.
\end{block}
\renewcommand{\arraystretch}{1.15}
\centering\footnotesize
\begin{tabular}{@{}>{\raggedright\arraybackslash}p{6.2cm}>{\raggedright\arraybackslash}p{6.2cm}@{}}
\toprule
\textbf{Verificación de estado} (balanceador) & \textbf{Verificación de tiempo de actividad} (Monitoring)\\
\midrule
Decide sola a qué VM enviar tráfico & Avisa a una persona que algo no responde\\
\bottomrule
\end{tabular}
\end{frame}

% 23
\begin{frame}{Tipos de balanceador}
\renewcommand{\arraystretch}{1.35}
\centering\small
\begin{tabular}{@{}p{2.6cm}>{\raggedright\arraybackslash}p{5.0cm}>{\raggedright\arraybackslash}p{5.0cm}@{}}
\toprule
 & \textbf{De aplicaciones} & \textbf{De red}\\
\midrule
\textbf{Protocolo} & HTTP y HTTPS & TCP, UDP y TCP con SSL\\
\textbf{Qué mira} & La URL de cada solicitud & Solo la IP y el puerto\\
\textbf{Qué puede hacer} & Enviar \texttt{/api} a un grupo de VMs y \texttt{/imagenes} a un bucket; manejar el certificado HTTPS & Repartir conexiones sin abrir su contenido\\
\textbf{Ejemplo} & La página y la app de Pedidos Express & Un servidor de juegos o de video en vivo\\
\bottomrule
\end{tabular}
\par\vspace{3mm}
\begin{columns}[T,onlytextwidth]
\column{.48\textwidth}
\begin{block}{Externo}
\small Recibe tráfico de Internet. Es el que ven los clientes.
\end{block}
\column{.48\textwidth}
\begin{block}{Interno}
\small Reparte tráfico entre servicios dentro de tu red VPC. No tiene IP pública.
\end{block}
\end{columns}
\end{frame}

% 24
\begin{frame}{Autoescalado: más VMs cuando hace falta}
\begin{columns}[c,onlytextwidth]
\column{.50\textwidth}
\begin{block}{Autoescalado}
\small El grupo de instancias agrega o quita VMs según una métrica. Usa la misma métrica que ves en Cloud Monitoring.
\end{block}
\begin{exampleblock}{Configuración de Pedidos Express}
\small Mínimo \textbf{2} VMs, máximo \textbf{6}.\\
Objetivo: uso de CPU promedio de \textbf{60\,\%}.
\end{exampleblock}
\column{.46\textwidth}
\renewcommand{\arraystretch}{1.3}
\centering\small
\begin{tabular}{@{}lrc@{}}
\toprule
\textbf{Hora} & \textbf{CPU promedio} & \textbf{VMs}\\
\midrule
04:00 & 10\,\% & 2 {\color{utpGris}(mínimo)}\\
13:00 & 75\,\% $\rightarrow$ 58\,\% & 2 $\rightarrow$ 3\\
\rowcolor{utpRojo!20} 20:00 & 90\,\% $\rightarrow$ 60\,\% & 3 $\rightarrow$ 5\\
23:30 & 20\,\% & 5 $\rightarrow$ 2\\
\bottomrule
\end{tabular}
\end{columns}
\par\vspace{4mm}
\begin{alertblock}{Cómo encaja con el balanceador}
\small El balanceador envía tráfico al grupo. Cuando el grupo crea una VM nueva, el balanceador empieza a usarla apenas pasa la verificación de estado.
\end{alertblock}
\end{frame}

% 25
\begin{frame}{Cloud CDN: copias cerca del cliente}
\begin{center}
\begin{tikzpicture}[font=\small]
  \node[caja,fill=utpRojo!12,minimum width=1.9cm,minimum height=0.9cm] (u) at (0,0) {cliente\\\footnotesize en Lima};
  \node[caja,fill=utpRojo!30,text width=2.6cm,minimum height=1.3cm] (cdn) at (4.3,0) {\textbf{Caché de Cloud CDN}\\\footnotesize cerca del cliente};
  \node[caja,fill=white,text width=2.4cm,minimum height=1.3cm] (lb) at (9.4,0) {\textbf{Balanceador}\\\footnotesize y VMs en \texttt{us-central1}};
  \draw[flecha] (u.east) -- node[above,font=\scriptsize]{\texttt{logo.png}} (cdn.west);
  \draw[flecha,dashed] (cdn.east) -- node[above,font=\scriptsize]{solo la 1.\textsuperscript{a} vez} (lb.west);
\end{tikzpicture}
\end{center}
\begin{columns}[T,onlytextwidth]
\column{.48\textwidth}
\begin{block}{Primera solicitud}
\small La caché no tiene \texttt{logo.png}: lo pide a las VMs, se lo entrega al cliente y \textbf{guarda una copia}.
\end{block}
\column{.48\textwidth}
\begin{block}{Siguientes solicitudes}
\small La caché responde con su copia. La solicitud \textbf{no llega} a las VMs.
\end{block}
\end{columns}
\par\vspace{2mm}
\renewcommand{\arraystretch}{1.15}
\centering\footnotesize
\begin{tabular}{@{}>{\raggedright\arraybackslash}p{6.2cm}>{\raggedright\arraybackslash}p{6.2cm}@{}}
\toprule
\textbf{Sí va en la CDN} (igual para todos) & \textbf{No va en la CDN} (cambia por cliente)\\
\midrule
Fotos de productos, logo, CSS, JavaScript & Carrito, pedidos de Ana, estado del pago\\
\bottomrule
\end{tabular}
\par\vspace{2mm}
\raggedright\small Cloud CDN se activa con una casilla en el balanceador de aplicaciones externo.
\end{frame}

% 26
\begin{frame}{Todo junto: Pedidos Express en línea}
\begin{center}
\begin{tikzpicture}[font=\small]
  \node[caja,fill=utpRojo!12,minimum width=1.6cm] (u) at (0,1.0) {clientes};
  \node[caja,fill=utpRojo!30,text width=2.5cm,minimum height=1.1cm] (lb) at (3.1,1.0) {\textbf{Balanceador}\\\footnotesize + Cloud CDN};
  \draw[flecha] (u) -- (lb);
  \draw[utpRojo,thick,rounded corners=4pt] (5.3,-1.25) rectangle (12.95,2.05);
  \node[font=\scriptsize\bfseries,color=utpRojo,anchor=south east] at (12.95,2.05) {región us-central1};
  \draw[thick] (lb.east) -- (5.75,1.0) |- (11.65,1.65);
  \foreach \z/\n/\x in {a/1/6.65,b/2/9.15,c/3/11.65} {
    \draw[zona] (\x-1.15,-1.1) rectangle (\x+1.15,1.9);
    \node[font=\scriptsize,color=utpGris,anchor=south] at (\x,-1.1) {zona \z};
    \node[vm] (w\n) at (\x,1.0) {\texttt{web-\n}};
    \draw[flecha] (\x,1.65) -- (w\n.north);
  }
  \node[caja,fill=utpRojo!10,font=\scriptsize,minimum width=1.9cm] (s1) at (6.65,-0.25) {Cloud SQL\\principal};
  \node[caja,fill=white,font=\scriptsize,minimum width=1.9cm] (s2) at (9.15,-0.25) {Cloud SQL\\en espera};
  \draw[-latex,dashed] (s1) -- node[above,font=\tiny]{copia} (s2);
  \node[caja,fill=white,text width=2.5cm,font=\footnotesize] (mon) at (3.1,-0.4) {\textbf{Cloud Monitoring}\\alertas por correo};
  \draw[dashed,utpGrisMedio,thick] (mon.east) -- (5.3,-0.4);
\end{tikzpicture}
\end{center}
\renewcommand{\arraystretch}{1.1}
\centering\footnotesize
\begin{tabular}{@{}>{\raggedright\arraybackslash}p{3.4cm}>{\raggedright\arraybackslash}p{9.2cm}@{}}
\toprule
\textbf{Si pasa esto\dots} & \textbf{\dots responde esto}\\
\midrule
Falla \texttt{web-1} & La verificación de estado la saca; \texttt{web-2} y \texttt{web-3} atienden\\
Llega la hora de la cena & El autoescalado agrega VMs; la CDN entrega las fotos sin tocar las VMs\\
Falla la base principal & Cloud SQL pasa a la copia en espera de la zona \texttt{b}\\
Algo sale de lo normal & Cloud Monitoring envía un correo al equipo\\
\bottomrule
\end{tabular}
\end{frame}

% 27
\begin{frame}[fragile]{Práctica 1: balanceo con Nginx en tu computadora}
\begin{block}{Caso de estudio}
\small Pedidos Express prepara su tienda online para el Cyber Day: tres copias del servidor (Node.js) detrás de \textbf{Nginx}, que reparte por turnos: 3001, 3002, 3003, 3001\dots{} (\textit{round robin}).
\end{block}
\begin{columns}[c,onlytextwidth]
\column{.38\textwidth}
\begin{tikzpicture}[font=\small]
  \node[caja,fill=utpRojo!12,minimum width=1.3cm,font=\footnotesize] (c) at (0,0) {clientes};
  \node[caja,fill=utpRojo!30,minimum width=1.3cm,align=center] (n) at (2.0,0) {\textbf{Nginx}\\\footnotesize :8080};
  \draw[flecha] (c) -- (n);
  \foreach \p/\y in {3001/0.9,3002/0,3003/-0.9} {
    \node[vm,minimum width=1.1cm,minimum height=0.6cm] (s\p) at (4.0,\y) {\texttt{\p}};
    \draw[flecha] (n.east) -- (s\p.west);
  }
\end{tikzpicture}
\column{.60\textwidth}
\renewcommand{\arraystretch}{1.15}
\scriptsize
\begin{tabular}{@{}lll@{}}
\toprule
\textbf{Ruta} & \textbf{Qué hace} & \textbf{El servidor\dots}\\
\midrule
\texttt{GET /} & Página principal & responde al instante\\
\texttt{GET /productos} & Consulta la base de datos & \textbf{espera} 2 s\\
\texttt{POST /comprar} & Cobra con el banco & \textbf{espera} 5 s\\
\texttt{GET /reporte} & Calcula un reporte & \textbf{trabaja} 5 s con la CPU\\
\bottomrule
\end{tabular}
\end{columns}
\par\vspace{3mm}
\begin{exampleblock}{Cómo empezar (carpeta \texttt{Sesion25/practica\_balanceo/})}
\footnotesize
\begin{tabular}{@{}l>{\ttfamily}l@{}}
1. Tres terminales, un servidor en cada una: & node server.js 3001 \quad\textsf{(y 3002, 3003)}\\
2. Otra terminal, para encender Nginx: & nginx -p \string"\$PWD\string" -c nginx.conf\\
3. Otra terminal, para simular clientes: & node prueba.js / 6 --en-fila\\
\end{tabular}
\end{exampleblock}
\end{frame}

% 28
\begin{frame}{Práctica 1: seis pruebas}
\renewcommand{\arraystretch}{1.12}
\centering\small
\begin{tabular}{@{}c>{\ttfamily\scriptsize\raggedright\arraybackslash}p{5.5cm}>{\raggedright\arraybackslash}p{6.6cm}@{}}
\toprule
\textbf{\#} & \textsf{\small\textbf{Comando}} & \textbf{Qué observar}\\
\midrule
1 & node prueba.js / 6 --en-fila & ¿En qué orden responden los servidores?\\
2 & node prueba.js /productos 6 --directo 3001 & Un solo servidor: ¿cuánto tarda en total?\\
3 & node prueba.js /reporte 6 --directo 3001 & Un solo servidor: ¿cuánto espera el sexto cliente?\\
4 & node prueba.js /reporte 6 & Con Nginx: ¿cuánto tarda en total?\\
5 & node prueba.js /reporte 1\newline \textsf{\footnotesize y, antes de 5 s, en otra terminal:}\newline node prueba.js / 6 & ¿Qué visitas tardan? ¿A qué servidor le tocaron?\\
6 & \textsf{\footnotesize Apaga el 3002 (}Ctrl+C\textsf{\footnotesize ) y luego:}\newline node prueba.js / 6 --en-fila & ¿Algún cliente ve un error? ¿Qué servidores responden?\\
\bottomrule
\end{tabular}
\par\vspace{3mm}
\begin{alertblock}{La idea}
\small Cuando el servidor \textbf{espera}, atiende a muchos a la vez. Cuando \textbf{trabaja} con la CPU, los clientes hacen cola: ahí ayuda repartir entre varios servidores.
\end{alertblock}
\end{frame}

% 29
\begin{frame}[c]{Práctica 2}
\begin{block}{Realiza el laboratorio}
\large \textit{Supervisa recursos con Google Cloud Observability}\par\vspace{1mm}
\small\texttt{skills.google/catalog\_lab/1785}
\end{block}
\begin{exampleblock}{Al terminar}
\small Tienes un panel con el uso de CPU de las VMs \texttt{nginxstack}, una política de alertas con dos condiciones, el grupo \texttt{Instancias de VM} y una verificación de tiempo de actividad.
\end{exampleblock}
\begin{alertblock}{No olvides}
\small Desactivar la política de alertas antes de finalizar el lab (tarea 6).
\end{alertblock}
\end{frame}

% 30
\begin{frame}{Cierre: preguntas para comprobar}
\begin{block}{Responde con tus palabras}
\small
1. ¿Qué diferencia hay entre una métrica y un registro? Da un ejemplo de cada uno.\\\medskip
2. La CPU de una VM está al 100\,\% a las 4 a.\,m., sin clientes. ¿Qué podría estar pasando y dónde lo revisas?\\\medskip
3. ¿Qué es un punto único de falla? ¿Cómo lo quitas en los servidores web de Pedidos Express?\\\medskip
4. ¿Qué hace el balanceador cuando una VM no pasa la verificación de estado?\\\medskip
5. ¿Qué archivos de Pedidos Express conviene entregar con Cloud CDN y cuáles no? ¿Por qué?
\end{block}
\begin{alertblock}{Idea final}
\small Supervisar te dice qué pasa. La alta disponibilidad y el balanceador hacen que una falla no deje a los clientes sin servicio.
\end{alertblock}
\end{frame}

% 31
\begin{frame}{Referencias}
\scriptsize
\begin{itemize}\setlength\itemsep{4pt}
  \item Laboratorio: \textit{Supervisa recursos con Google Cloud Observability} (\textit{Monitor Resources with Google Cloud Observability}), curso \textit{Essential Google Cloud Infrastructure: Core Services}: \texttt{skills.google/catalog\_lab/1785}.
  \item Revisión: \texttt{github.com/rgap/rgap-Integracion-de-Servicios-Cloud-Gcp-Agosto2026}, carpeta \texttt{Sesion25/practica/} (\texttt{practica\_bigquery\_python.ipynb} y \texttt{generar\_ventas.py}). DuckDB: \texttt{duckdb.org}.
  \item Práctica 1: mismo repositorio, carpeta \texttt{Sesion25/practica\_balanceo/} (\texttt{README.md}, \texttt{server.js}, \texttt{nginx.conf} y \texttt{prueba.js}). Nginx: \texttt{nginx.org}. Node.js: \texttt{nodejs.org}.
  \item Material previo del curso: Sesión 24, Big Data y BigQuery.
  \item Google Cloud: Cloud Monitoring, Cloud Logging, Cloud Load Balancing, Cloud CDN, grupos de instancias administrados y Cloud SQL.
  \item Las respuestas de las preguntas están en \texttt{../respuestas.md}.
\end{itemize}
\end{frame}

\end{document}
